\section{Skeletons and constructive periods}
\label{sec:skeletons}

The affine term first gives periodicity at each coordinate, with an
argument that remains valid for composite bases.

\begin{lemma}[Coordinate periodicity]
\label{lem:coordinate-periodicity}
  Let $k\in\Z$ and $e=\nu_p(L_p(k))$.  Then $p^{e+1}$ is a period of the
  position $k$ in $x_{p,u}$.
\end{lemma}

\begin{proof}
  Write $L_p(k)=p^e a$ with $p\nmid a$.  For every $t\in\Z$,
  \begin{equation}
    L_p\bigl(k+t p^{e+1}\bigr)
      =p^e\bigl[a+(p-1)tp\bigr].
    \label{eq:coordinate-periodicity}
  \end{equation}
  The bracket is congruent to $a$ modulo $p$, so it is not divisible by
  $p$.  The exponent of \cref{eq:coordinate-periodicity} is therefore
  $e$, and the directive letter is unchanged along the whole arithmetic
  progression.  No inverse modulo $p$ was used.
\end{proof}

The centers in \cref{eq:L-r} satisfy two identities that will recur:
\begin{equation}
  L_p(r_N)=p^N,
  \qquad
  r_{N+1}=r_N+p^N.
  \label{eq:center-identities}
\end{equation}

\begin{theorem}[Exact one-hole skeleton]
\label{thm:exact-skeleton}
  For every $N\geq1$,
  \[
    \Per_{p^N}(x_{p,u})=\Z\setminus(r_N+p^N\Z).
  \]
\end{theorem}

\begin{proof}
  Since $\gcd(p-1,p)=1$, the congruence in \cref{eq:L-r} gives
  \begin{equation}
    p^N\mid L_p(k)
      \quad\Longleftrightarrow\quad
    k\equiv r_N\pmod {p^N}.
    \label{eq:hole-congruence}
  \end{equation}
  If $k$ is outside this residue class, put
  $e=\nu_p(L_p(k))<N$ and write $L_p(k)=p^e a$, where $p\nmid a$.  Then,
  for every $t\in\Z$,
  \begin{equation}
    L_p(k+t p^N)
       =p^e\bigl[a+(p-1)t p^{N-e}\bigr].
    \label{eq:skeleton-nonhole}
  \end{equation}
  The bracket in \cref{eq:skeleton-nonhole} is $a$ modulo $p$; hence the
  exponent and letter remain constant.  Thus every nonhole position lies
  in $\Per_{p^N}(x_{p,u})$.

  Conversely, take any $k\equiv r_N\pmod {p^N}$.  Its $p^N$-progression
  contains both $r_N$ and $r_{N+1}=r_N+p^N$.  By
  \cref{eq:center-identities}, their letters are $u_N$ and $u_{N+1}$.
  They are unequal by \cref{eq:neighbor-distinct}, so no coordinate in the
  alleged hole progression is $p^N$-periodic.  This proves the reverse
  inclusion and the equality.
\end{proof}

\begin{corollary}[Essential periods and simple Toeplitz form]
\label{cor:simple-toeplitz}
  Every $p^N$ is essential, the powers $(p^N)_{N\geq1}$ form a period
  structure, and $x_{p,u}$ is an aperiodic normal simple Toeplitz sequence.
\end{corollary}

\begin{proof}
  The $p^N$-skeleton has exactly one hole residue modulo $p^N$.  Any
  positive translation period of that partial word must preserve the hole
  set, and hence must be divisible by $p^N$.  Its least period is therefore
  $p^N$.

  Define
  \begin{equation}
    H_N=\{k:p^N\mid L_p(k)\}=r_N+p^N\Z,
    \qquad H_0=\Z.
    \label{eq:nested-holes}
  \end{equation}
  The layer $H_N\setminus H_{N+1}$ consists exactly of coordinates with
  exponent $N$, so it is filled uniformly by $u_N$, leaving the one hole
  class $H_{N+1}$.  No integer lies in every $H_N$: otherwise the fixed
  nonzero integer $L_p(k)$ would be divisible by all powers of $p$.
  Together with \cref{lem:coordinate-periodicity}, this shows that the
  periodic positions cover $\Z$ and that the one-hole filling has no
  residual coordinate.  This is the normal simple Toeplitz recursion.

  If $x_{p,u}$ had a nonzero global period $q$, translation by $q$ would
  preserve each exact skeleton and therefore each unique hole class in
  \cref{eq:nested-holes}.  Hence $p^N\mid q$ for every $N$, impossible for
  a fixed nonzero integer.  Thus the sequence is aperiodic.
\end{proof}

Skeleton essentiality does not by itself settle constructiveness: the
latter concerns the common period of all positions in a finite block.  The
next theorem separates these notions.

\begin{theorem}[Constructiveness is equivalent to primality]
\label{thm:constructive-split}
  Let $B_N=x_{p,u}[0,p^N-1]$.  Its essential period equals $p^{N+1}$ for
  every $N\geq1$ if $p$ is prime.  If $p$ is composite and $\ell$ is any
  prime divisor of $p$, then $\ell p^N$ is a common period of $B_N$ and
  satisfies $\ell p^N<p^{N+1}$.
\end{theorem}

\begin{proof}
  We first obtain a common upper period for every integer base.  If
  $0\leq k<p^N$, then
  \[
    0<L_p(k)<p^{N+1}.
  \]
  Hence $e=\nu_p(L_p(k))\leq N$.  The calculation in
  \cref{lem:coordinate-periodicity} shows that $p^{N+1}$ is a period of
  every such position.  Thus $p^{N+1}$ is a common period of $B_N$.

  Suppose now that $p$ is prime and that a positive common period $q$ is
  not divisible by $p^{N+1}$.  Write
  \begin{equation}
    q=p^j d,\qquad 0\leq j\leq N,\qquad p\nmid d.
    \label{eq:prime-period-decomposition}
  \end{equation}
  Because $(p-1)d$ is invertible modulo $p^2$, there is $t\in\Z$ with
  \begin{equation}
    1+(p-1)dt\equiv p\pmod {p^2}.
    \label{eq:prime-witness}
  \end{equation}
  The coordinate $r_j$ lies in $[0,p^N-1]$, where $r_0=0$, and
  \begin{align*}
    L_p(r_j)&=p^j,\\
    L_p(r_j+tq)&=p^j\bigl[1+(p-1)dt\bigr].
  \end{align*}
  By \cref{eq:prime-witness}, the two divisibility exponents are exactly
  $j$ and $j+1$.  Their letters $u_j$ and $u_{j+1}$ differ, contradicting
  that $q$ is a period of the position $r_j$.  Thus every common period is
  divisible by $p^{N+1}$, and the upper bound gives equality.

  Finally let $p$ be composite, fix a prime $\ell\mid p$, and put
  $q=\ell p^N$.  Since $\ell<p$, this period is strictly smaller than
  $p^{N+1}$.  Fix $0\leq k<p^N$, let
  $e=\nu_p(L_p(k))$, and take an arbitrary $t\in\Z$.  If $e<N$, write
  $L_p(k)=p^e a$ with $p\nmid a$.  Then
  \begin{equation}
    L_p(k+tq)
       =p^e\bigl[a+(p-1)t\ell p^{N-e}\bigr].
    \label{eq:composite-low}
  \end{equation}
  The bracket is $a$ modulo $p$, so the exponent remains $e$.  If $e=N$,
  \cref{eq:hole-congruence} and $0\leq k<p^N$ force $k=r_N$, and
  \begin{equation}
    L_p(k+tq)=p^N\bigl[1+(p-1)\ell t\bigr].
    \label{eq:composite-top}
  \end{equation}
  The bracket is $1$ modulo $\ell$, so it cannot be divisible by $p$.
  Again the exponent is unchanged.  Since $t$ was arbitrary,
  $\ell p^N$ is a full-progression period of every position of $B_N$.
\end{proof}

\begin{figure}[!b]
  \centering
  \resizebox{\textwidth}{!}{% Reproducible monochrome TikZ source for Figure 2.
\begin{tikzpicture}[
  font=\small,
  >=Stealth,
  lane/.style={draw=black, rounded corners=2pt, inner sep=7pt,
               text width=.41\linewidth, align=left},
  proofbox/.style={draw=black, rounded corners=1.2pt, inner sep=5pt,
               text width=.36\linewidth, align=center},
  flow/.style={->, semithick}
]
  \node[lane,anchor=north west] (prime) at (0,0) {\textbf{Prime base $p$.}\\
    Let $q$ be a common period with $p^{N+1}\nmid q$.  Write
    $q=p^j d$, where $0\leq j\leq N$ and $p\nmid d$.};
  \node[proofbox,below=5mm of prime] (cong) {Choose $t\in\mathbb Z$ so that\\
    $1+(p-1)dt\equiv p\pmod {p^2}$.};
  \node[proofbox,below=5mm of cong] (kill) {$r_j\in[0,p^N-1]$ and\\
    $\nu_p(L_p(r_j))=j$,\quad
    $\nu_p(L_p(r_j+tq))=j+1$.};
  \node[proofbox,below=5mm of kill,fill=black!8] (presult)
    {$u_j\neq u_{j+1}$ contradicts common periodicity;\\
     the least common period is $p^{N+1}$.};
  \draw[flow] (prime) -- (cong);
  \draw[flow] (cong) -- (kill);
  \draw[flow] (kill) -- (presult);

  \node[lane,anchor=north west] (comp) at (.52\linewidth,0)
    {\textbf{Composite base $p$.}\\
     Choose a prime $\ell\mid p$ and set
     $q=\ell p^N<p^{N+1}$.  Fix $0\leq k<p^N$ and every $t\in\mathbb Z$.};
  \node[proofbox,below=5mm of comp] (low) {If $e=\nu_p(L_p(k))<N$, then\\
    $L_p(k+tq)=p^e[a+(p-1)t\ell p^{N-e}]$,\\
    whose bracket is $a\pmod p$.};
  \node[proofbox,below=5mm of low] (top) {If $e=N$, then $k=r_N$ and\\
    $L_p(k+tq)=p^N[1+(p-1)\ell t]$,\\
    whose bracket is $1\pmod\ell$.};
  \node[proofbox,below=5mm of top,fill=black!8] (cresult)
    {Every exponent is preserved for every translate;\\
     $\ell p^N$ is a strict common period.};
  \draw[flow] (comp) -- (low);
  \draw[flow] (low) -- (top);
  \draw[flow] (top) -- (cresult);
\end{tikzpicture}
}
  \caption{Schematic proof map for the finite-block dichotomy.  In the
  prime lane, every candidate common period missing a factor of
  $p^{N+1}$ is contradicted at a center $r_j$ by a congruence modulo
  $p^2$.  In the composite lane, each prime divisor $\ell$ of $p$ supplies
  the strict common period $\ell p^N$; the exponent-$N$ coordinate is
  checked separately.  The theorem quantifies over every integer
  translate, not only the displayed witnesses.}
  \label{fig:constructive-split}
\end{figure}

The proof covers $p=3$ and every composite base without treating a
nonmultiple of $p$ as a unit outside the prime lane.  If one adopts a
literal finite-block convention omitting coordinate zero, the same split
holds; the replacement witness is recorded in
\cref{app:boundary-conventions}.
